% mode: ini
% Aimed at write_tounicode's first lookup of a glyph, tfm:<tfmname>/<glyph>, which is left out when no name of either table
% begins with that prefix: names that sort next to tfm:cmr10/ without beginning with it, a font whose prefix some names have,
% a name that is the prefix alone, and the components of glyph names with `_`, each looked up under the prefix.
\input prelude
\pdfoutput=1 \pdfcompresslevel=0 \pdfobjcompresslevel=0
\pdfhorigin=1in \pdfvorigin=1in \pdfpagewidth=300pt \pdfpageheight=200pt

\pdfmapline{x}
\pdfmapline{+cmr10 CMR10 <cmr10.pfb}
\pdfmapline{+cmr5 CMR5 <cmr5.pfb}
\pdfmapline{+cmr7 CMR7 <pdffont/tfmprefix.enc <cmr7.pfb}
\pdfmapline{+cmr9 CMR9 <cmr9.pfb}
\pdfglyphtounicode{a}{0061}
\pdfglyphtounicode{b}{0062}
\pdfglyphtounicode{i}{0069}
\pdfglyphtounicode{tfm:cmr/a}{0030}
\pdfglyphtounicode{tfm:cmr1/a}{0031}
\pdfglyphtounicode{tfm:cmr10x/a}{0032}
\pdfglyphtounicode{tfm:cmr10-/b}{0033}
\pdfglyphtounicode{tfm:cmr100/a}{0034}
\pdfglyphtounicode{tfm:cmr10.5/a}{0035}
\pdfglyphtounicode{tfm:cmr5/a}{0036}
\pdfglyphtounicode{tfm:cmr5/b}{0037 0038}
\pdfglyphtounicode{tfm:cmr5/c}{xyz}
\pdfglyphtounicode{tfm:cmr9/}{0039}
\pdfglyphtounicode{tfm:cmr7/f}{0046}
\pdfglyphtounicode{tfm:cmr7/a}{0041}
\pdfgentounicode=1
\font\a=cmr10 \font\b=cmr5 \font\c=cmr7 \font\d=cmr9
\setbox0=\hbox{\a abc\b abc\c abcdefghi\d abc}
\shipout\box0
\end
